\section{Backgrounds}\label{Backgrounds}
  Each character has a history from before the current campaign.
  This can include their childhood, family, previous occupations, and more.
  Backgrounds generally do not have direct effects on a character's statistics.
  Narratively, a background explains a character's statistics, not defines them.
  However, a character's background can still have a significant influence on how they interact with the world.

  Backgrounds can be summarized as a set of benefits and flaws.
  If you choose a background benefit, you must also choose a flaw.
  Some backgrounds internally provide both benefits and flaws.
  These are called mixed backgrounds, and taking them does not change how many benefit or flaw backgrounds you can choose.
  You generally shouldn't have more than one benefit and one flaw, or a single mixed background.

  Not every character should have a specific background with effects listed here.
  These backgrounds can significantly change how a character relates to the world.
  It's fine to have a simpler background to put more focus on other aspects of your character's narrative journey.

  At the GM's discretion, background benefits can also be acquired during a campaign.
  For example, successfully performing a heroic and extremely public feat might give the whole party the benefits of the Folk Hero background.
  Conversely, a party who is defeated in battle may find themselves with the Indebted background, as the victors decide to spare them - with a price.

  Many backgrounds, especially background benefits, are only relevant in some area where your reputation could plausibly be known.
  If you're deep in the tunnels of Terra, nobody will care that you were a folk hero back home.
  The GM should ensure that backgrounds are usually relevant, and your reputation may become more broadly known through the course of a campaign.

  \subsection{Background Benefits}
    \parhead{Apothecary's Assistant} You spent years working with an herbalist or apothecary.
    You have a knack for identifying plants, herbs, and fungi in the wild.
    In addition, you can often find local suppliers of rare ingredients who might otherwise be wary of strangers.
    \parhead{Criminal Connections} You have a reputation in criminal circles as a trustworthy partner in crime.
    This reputation is not generally known outside of that domain, so law enforcement will not generally cause you problems.
    You know how to reliably identify and contact other criminals, at least in a significant local area.
    They will generally be helpful, though they may still charge a fair price for any direct assistance.
    \parhead{Folk Hero} You are known generally by most common folk as a heroic and benevolent figure, at least in a significant local area.
    This reputation could be well-earned or built on deception.
    In either case, it is widely believed.
    Common people will generally try to be helpful whenever you need it, as long as you act appropriately to match the tales.
    \parhead{Guild Member} You are a member of a major trade guild.
    This could be based on your abilities with a particular Craft skill, your class, or some other abilities you have that are common enough in the world to merit a trade guild.
    The GM can provide more guidance about what guilds exist.
    In general, the more advanced and civilized the world is, the more niche guilds exist.
    Your guild will have a presence in any major city and some minor ones.
    People will generally be more willing to believe that you can perform tasks relevant to your guild membership.
    In addition, fellow guild members will be more helpful when you need it.
    \parhead{Mysterious Heirloom} You have a family heirloom that you must keep safe.
    In the right hands, it would probably have great value or importance, but most people would dismiss it as a simple trinket or oddity.
    \parhead{Landed} You have a family estate that you can live in comfortably, including people who work the land and maintain the house.
    It is well equipped for most types of training, research, or other long-term work you and your allies might need to do.
    You do not have sole ownership of the estate, so selling it or otherwise attempting to directly gain value from it would be complicated and risky.
    However, it is a safe refuge when you need it.
    \parhead{Noble} Your family is at least minor nobility.
    This can change people's reactions to you and allow you to more easily access events and important people in high society.
    \parhead{Veteran} You have served as a low-level soldier of some type in a formal military organization.
    Other soldiers and veterans will generally treat you with a degree of respect and camaraderie.
    You have easier access to local military officers, and they will be more willing to help you out when you need it.

  \subsection{Background Flaws}
    \parhead{Cursed Legacy} Your family bears a persistent curse that has manifested in you.
    It is not life-threatening, but it causes you hardship, such as bad luck or periodic bouts of pain.
    Most commoners who find out about your curse will react with suspicion or fear.
    The curse may be removed, but not by means as simple as the \ritual{dispel curse} ritual.
    \parhead{Escapee} You have escaped from some private individual or institution who is very interested in your return.
    This could be because you were a slave, because of something unique about your heritage, or any other reason.
    The people who want your return will invest time and resources into pursuing you wherever it is feasible, including hiring bounty hunters to track you down and bring you back.
    % \parhead{Forsworn} You are generally known to be a liar who has broken an important oath.
    % Anyone aware of this fact will not trust you or your apparent allies with anything of importance.
    \parhead{Indebted} Thanks to your own past mistakes or your family's actions, you owe a significant debt.
    Directly paying off the debt would require at least a rank 3 wealth item (1,000 gp).
    The creditor may require you to make smaller recurring payments, or may compel you to perform tasks for them to pay down your debt over time.
    They can generally be negotiated with to a limited degree, but if the relationship turns sour, this flaw may switch to the Escapee flaw.
    \parhead{Nemesis} Someone specifically wants you to suffer due to your shared history.
    They and their allies are initially stronger than you in a direct conflict, and directly attacking them is almost certainly illegal.
    You will have to avoid directly confronting them, at least at first.
    Unlike the Escapee and Wanted flaws, your nemesis will not generally try to kill you or physically harm you.
    Instead, they will act to subvert your goals, either directly or through agents.
    Anything you seem to want, they will do their best to thwart.
    \parhead{Oathbound} You have sworn a binding vow to a deity, a powerful individual, or a secret society.
    This oath places significant restrictions on your behavior, such as a vow of silence in certain contexts, a requirement to destroy certain types of items, or a prohibition against killing certain types of creatures.
    Breaking this oath would have severe social or personal consequences, including but not limited to enmity from whatever entity you swore the oath to.
    \parhead{Repulsive} You are personally noxious, odorous, grotesque, or otherwise unpleasant to see and spend time with.
    The cause could be injury, disease, a powerful curse, or some other reason.
    This negatively affects almost everyone's reactions to you, which makes social interactions more difficult.
    In addition, any living creature who can see, hear, or smell you is unable to benefit from a \glossterm{short rest} or \glossterm{long rest}.
    Even if you can find loyal travelling companions, they still have their limits, forcing you to camp alone.
    \parhead{Wanted} You are wanted for serious crimes in some major area.
    This flaw does not require that there are posters with your name in every city, but the area where you are wanted must be important to the campaign.
    You will have to avoid identification by both law enforcement and even common people while in that area, and in other areas that may be aware of those details.
    In addition, bounty hunters may pursue you wherever you go.
    You may be innocent of the charges, but if so, it would be hard to prove your innocence.

  \subsection{Mixed Backgrounds}
    \parhead{Scion} You are in the line of inheritance for an important throne or noble house.
    You are not the designated heir, but some small distance removed, such as a third child.
    If the obstacles to your inheritance were cleared, you could become powerful and wealthy.
    However, you may also be a target for people trying to ensure their own inheritance, or simply using you to manipulate your family.
    As an important figure, your family or related people may try to place restrictions on your actions to ensure your safety.

    % \section{Sample Characters}

    %     This section lists sample characters for each class archetype.
    %     You can simply pick up one of these characters and use it as your character.
    %     Alternately, you can use a sample character as a starting point and adjust it to match your own character concept.
    %     The sample characters are ordered by class first, and by archetype within each class second.

    %     \subsection{Barbarian}

    %         \subsubsection{Battleforged Resilience}
    %             \parhead{Species} Dwarf.
    %             % 2 1 3 0 2 2 base
    %             \parhead{Attributes} 2 Str, -1 Dex, 4 Con, 0 Int, 2 Per, 2 Wil (after species modifiers).
    %             \parhead{Class} Barbarian.
    %             \parhead{Archetypes} Battleforged Resilience first, Primal Warrior second, Totemist (bear totem) third.
    %             \parhead{Insight Points} 1.
    %             \parhead{Skills} Awareness, Athletics, Endurance
    %             \parhead{Languages} Common, Dwarven, Giantish.
    %             \parhead{Equipment} Battleaxe, standard shield, scale mail. As you gain levels, use the best medium armor you can afford. If you gain proficiency with heavy armor, use that instead.
    %             \parhead{Combat Styles} Herald of War, Unbreakable Defense.
    %             \parhead{Suggested Feats} Shieldbearer, Martial Training, Regenerator, Toughness.
    %             \parhead{Combat Tactics} You are extremely difficult to kill.
    %             Take advantage of that by wading into the front lines of combat and drawing attention away from your more vulnerable allies.
    %             If you find yourself in danger, use defensive maneuvers like \maneuver{defensive strike} and \maneuver{flamboyant parry} to keep yourself safe.
    %             On the other hand, if your foes try to ignore you after realizing how durable you are, force them to engage with you using maneuvers like \maneuver{challenging strike} and \maneuver{guard the pass}.

    %         \subsubsection{Battlerager}
    %             \parhead{Species} Half-orc.
    %             % 3 2 2 0 2 1 base
    %             \parhead{Attributes} 4 Str, 2 Dex, 2 Con, -1 Int, 1 Per, 0 Wil (after species modifiers).
    %             \parhead{Class} Barbarian.
    %             \parhead{Archetypes} Battlerager first, Primal Warrior second, Totemist (lion totem) third.
    %             \parhead{Insight Points} 0.
    %             \parhead{Skills} Awareness, Endurance, Intimidate.
    %             \parhead{Equipment} Greatmace, scale mail. As you gain levels, buy a heavy crossbow and use the best medium armor you can afford.
    %             \parhead{Combat Styles} Herald of War.
    %             \parhead{Suggested Feats} Greatweapon Warrior, Rapid Reaction, Swiftrunner.
    %             \parhead{Combat Tactics} You are a furious frenzy of devastating damage and lethal critical hits.
    %             When you roll a 10 on an attack roll, whatever you attacked will probably die.
    %             Staying close to your allies is generally a good plan, since you don't have the durability to run into the middle of a horde of enemies safely.
    %             Your maneuvers help you deal with high-Armor enemies and enemy swarms, and give you the ability to sacrifice most of your statistics other than damage in exchange for more damage.

    %         \subsubsection{Outland Savage}
    %             \parhead{Species} Elf.
    %             % 2 3 2 1 2 0 base
    %             \parhead{Attributes} 2 Str, 4 Dex, 1 Con, 1 Int, 1 Per, -1 Wil (after species modifiers).
    %             \parhead{Class} Barbarian.
    %             \parhead{Archetypes} Outland Savage first, Primal Warrior second, Totemist (wolf totem) third.
    %             \parhead{Insight Points} 2 points (1 for proficiency with exotic armor weapons, 1 for Mobile Hunter combat style).
    %             \parhead{Skills} Awareness, Athletics, Endurance, Stealth, Survival.
    %             \parhead{Languages} Common, Orcish.
    %             \parhead{Equipment} Flail, standard shield, scale mail. As you gain levels, use the best light armor you can afford. When you can, get spikes and a spiked knee crafted onto your armor.
    %             \parhead{Combat Styles} Dirty Fighting, Mobile Hunter
    %             \parhead{Suggested Feats} Savage, Brawler, Swiftrunner.
    %             \parhead{Combat Tactics} You can move around the battlefield very quickly, and you are incredibly accurate with special combat actions like shoving and grappling enemies.
    %             Make the most of that by repositioning enemies, tripping them, or holding them in grapples so your allies can hit them.
    %             While you aren't in a grapple, use your flail in two hands to maximize your damage.
    %             When you enter a grapple, use your spiked knee to attack, since your flail is much less effective while grappling.
    %             If you don't have any allies who like being on the front lines, you won't be as effective at helping them deal damage to enemies, but you're still very skilled at preventing enemies from reaching your allies.
    %             In that case, consider choosing bear totem or shark totem instead of wolf totem.

    %         \subsubsection{Primal Warrior}
    %             \parhead{Species} Human.
    %             \parhead{Attributes} 3 Str, 2 Dex, 2 Con, 1 Int, 2 Per, 0 Wil.
    %             \parhead{Class} Barbarian.
    %             \parhead{Archetypes} Primal Warrior first, Battleforged Resilience second, Outland Savage third.
    %             \parhead{Insight Points} 1 point for an additional combat style, 1 point for an additional maneuver.
    %             \parhead{Skills} Awareness, Athletics, Endurance, Intimidate.
    %             \parhead{Languages} Common, Dwarven, Orcish.
    %             \parhead{Equipment} Greataxe, scale mail. As you gain levels, buy a heavy crossbow and use the best medium armor you can afford.
    %             \parhead{Combat Styles} Dirty Fighting, Herald of War, Unbreakable Defense.
    %             \parhead{Suggested Feats} Greatweapon Warrior, Weapon Focus, Swiftrunner.
    %             \parhead{Combat Tactics} You have a great breadth of options available to you thanks to the number of maneuvers you know.
    %             You have the survivability to stand in close combat, especially if you use maneuvers from Unreakable Defense, but you can also shout at mobile enemies from range with maneuvers from Herald of War.
    %             Both Dirty Fighting and Herald of War give you maneuvers that work well against enemies with a high Armor defense, so you can adapt to whatever battle you find yourself in.
    %             You can make the most of your versatility by learning maneuvers like \maneuver{knockback shove} that are sometimes useless, but which can be devastatingly effective in the right context.

    %         \subsubsection{Totemist}
    %             Characters from this archetype can be very different based on their chosen totem.
    %             A bear totem character might resemble the typical character for the Battleforged Resilience archetype.
    %             A lion totem or shark totem character might resemble the typical character for the Battlerager archetype.
    %             A wolf totem character might resemble the typical character for the Outland Savage archetype.

    %             If you want to quickly create a character based on the eagle totem from this archetype, make the following choices:
    %             \parhead{Species} Human.
    %             \parhead{Attributes} 2 Str, 2 Dex, 3 Con, 0 Int, 4 Per, -1 Wil.
    %             \parhead{Class} Barbarian.
    %             \parhead{Archetypes} Totemist (eagle totem) first, Primal Warrior second, Outland Savage third.
    %             \parhead{Insight Points} 1 point for Perfect Precision combat style.
    %             \parhead{Skills} Awareness, Poise, Athletics, Creature Handling, Survival.
    %             \parhead{Languages} Common, Elven, Giantish.
    %             \parhead{Equipment} Longbow, leather body armor. As you gain levels, use the best light armor you can afford.
    %             \parhead{Combat Styles} Perfect Precision.
    %             \parhead{Suggested Feats} Sniper, Blindfighter, Swiftrunner.
    %             \parhead{Combat Tactics} You have incredible accuracy from very long range.
    %             Your defenses are relatively low, but as long as you stay far enough away from your foes, they can't take advantage of that weakness.
    %             You have the ability to prioritize any target on the battlefield, so make the most of your maneuvers that impose conditions or deal additional damage on weakened foes.

    %     \subsection{Cleric}

    %         \subsubsection{Divine Magic}
    %             \parhead{Species} Gnome.
    %             % 0 0 2 1 2 3 base
    %             \parhead{Attributes} -1 Str, 0 Dex, 3 Con, 1 Int, 2 Per, 4 Wil (after species modifiers).
    %             \parhead{Class} Cleric.
    %             \parhead{Archetypes} Divine Magic first, Divine Spell Mastery second, Domain Influence third.
    %             \parhead{Insight Points} 3 points (2 points for an additional mystic sphere (Photomancy), 1 point for an additional rank 1 spell).
    %             \parhead{Skills} Knowledge (local, religion), Medicine, Persuasion, Social Insight
    %             \parhead{Languages} Common, Gnome, Halfling.
    %             \parhead{Equipment} Mace, standard shield, scale mail. As you gain levels, use the best medium armor you can afford.
    %             \parhead{Domains} Good, Magic
    %             \parhead{Mystic Spheres} Channel Divinity and Photomancy
    %             \parhead{Suggested Feats} Celestial Ancestry, Sphere Focus: Photomancy, Sphere Focus: Prayer
    %             \parhead{Combat Tactics} You can protect yourself and your allies and invoke divine wrath on your foes.
    %             Your attacks can hit a variety of defenses, so use the best spells for the situation.
    %             If you are facing a foe that not particularly vulnerable to your attacks, you can focus on helping your allies with ``boon'' spells to make their actions more effective and keep them safe.

    %         \subsubsection{Divine Spell Mastery}
    %             Use the typical character for the Divine Magic cleric archetype.
    %             Even if you focus on spells through this archetype, you should generally still rank up your spells before improving your rank in this archetype.

    %         \subsubsection{Domain Influence}
    %             Characters from this archetype can be very different based on their chosen domains.
    %             A character with spellcasting-focused domains might resemble the typical character for the Divine Magic cleric archetype.
    %             If you want to quickly create a more martial character based on the Strength and War domains from this archetype, make the following choices:

    %             \parhead{Species} Dwarf.
    %             % 2 0 3 0 2 1 base
    %             \parhead{Attributes} 2 Str, -1 Dex, 4 Con, 0 Int, 2 Per, 2 Wil (after species modifiers).
    %             \parhead{Class} Cleric.
    %             \parhead{Archetypes} Domain Influence first, Divine Magic second, Preacher third.
    %             \parhead{Insight Points} 2 points (unspent).
    %             \parhead{Skills} Awareness, Knowledge (local, religion), Medicine
    %             \parhead{Languages} Common, Draconic, Dwarven.
    %             \parhead{Equipment} Morning star, standard shield, scale mail. As you gain levels, use the best heavy armor you can afford.
    %             \parhead{Domains} Destruction, War
    %             \parhead{Mystic Sphere} Channel Divinity
    %             \parhead{Suggested Feats} Weapon Focus, Sphere Focus: Channel Divinity, Shieldbearer
    %             \parhead{Combat Tactics} You are a frontline fighter first and foremost.
    %             Your magically enhanced resistance and high defenses make you durable in combat, though you lack mobility. 
    %             When you need to distract foes or face down hordes, you can use your abilities from the Preacher archetype.

    %         \subsubsection{Healer}
    %             \parhead{Species} Human.
    %             % 0 1 2 1 0 2 base
    %             \parhead{Attributes} 0 Str, 2 Dex, 2 Con, 1 Int, 0 Per, 3 Wil.
    %             \parhead{Class} Cleric.
    %             \parhead{Archetypes} Healer first, Divine Magic second, Domain Influence third.
    %             \parhead{Insight Points} 3 points (2 points for an additional mystic sphere (Vivimancy), 1 point for an additional rank 1 spell).
    %             \parhead{Skills} Awareness, Deduction, Knowledge (local, religion), Medicine
    %             \parhead{Languages} Common, Draconic, Halfling.
    %             \parhead{Equipment} Morning star, standard shield, scale mail. As you gain levels, use the best medium armor you can afford.
    %             \parhead{Domains} Life, Protection
    %             \parhead{Mystic Spheres} Prayer, Vivimancy
    %             \parhead{Suggested Feats} Sphere Focus: Vivimancy, Boongiver, Sphere Focus: Prayer
    %             \parhead{Combat Tactics} You have an unmatched mastery of healing and protection.
    %             You have reasonably high defenses, so you can take to the front lines as necessary to make the most of your \ability{divine aid} and \ability{divine protection} abilities.
    %             Although your \ability{healer's grace} ability is powerful, you shouldn't feel bad about attacking enemies.
    %             That's especially important early in a fight when your allies don't need healing yet and your enemies haven't realized that it's pointless to attack your allies while you are still standing.

    %         \subsubsection{Preacher}
    %             \parhead{Species} Human.
    %             % 0 0 2 1 3 0 base
    %             \parhead{Attributes} 0 Str, 0 Dex, 2 Con, 2 Int, 4 Per, 0 Wil.
    %             \parhead{Class} Cleric.
    %             \parhead{Archetypes} Preacher first, Divine Magic second, Divine Spell Mastery third.
    %             \parhead{Insight Points} 4 points (2 points for an additional mystic sphere (Revelation), 2 points unspent).
    %             \parhead{Skills} Awareness, Knowledge (local, religion), Medicine, Persuasion, Social Insight
    %             \parhead{Languages} Common, Dwarven, Elven.
    %             \parhead{Equipment} Club, standard shield, scale mail. As you gain levels, use the best medium armor you can afford.
    %             \parhead{Mystic Spheres} Enchantment, Revelation
    %             \parhead{Suggested Feats} Persuasion Specialization, Sphere Focus: Enchantment, Sphere Focus: Revelation
    %             \parhead{Combat Tactics} Your social skills are virtually unmatched, and you have a wide variety of spells that give you narrative power in social situations.
    %             In combat, your \ability{denounce the heathens} ability is essentially guaranteed to hit, so you should stay close enough to the front lines to make good use of it.
    %             You can take advantage of the lowered defenses of your denounced foes to succeed with powerful mind-affecting spells.

    %     \subsection{Druid}

    %         \subsubsection{Elementalist}
    %             \parhead{Species} Human.
    %             % 0 0 2 1 3 2 base
    %             \parhead{Attributes} 0 Str, 0 Dex, 2 Con, 2 Int, 4 Per, 2 Wil.
    %             \parhead{Class} Druid.
    %             \parhead{Archetypes} Nature Magic first, Elementalist second, Nature Spell Mastery third.
    %             \parhead{Insight Points} 3 points (2 points for an additional mystic sphere, 1 point unspent).
    %             \parhead{Skills} Awareness, Poise, Endurance, Knowledge (dungeoneering, nature), Survival
    %             \parhead{Languages} Common, Sylvan
    %             \parhead{Equipment} Sickle, buckler, hide armor. As you gain levels, keep using hide armor.
    %             You may want to keep leather armor around in case you need to do a lot of jumping or swimming - or at high levels, flying.
    %             \parhead{Mystic Spheres} Any two of the four elemental mystic spheres.
    %             Your \textit{elemental spell} ability gives you access to spells from the fourth mystic sphere.
    %             That means that the specific three mystic spheres you choose mostly just affect which wands you can use and which feats you can take.
    %             \parhead{Suggested Feats} Sphere Focus: Aeromancy, Aquamancy, Pyromancy, or Terramancy
    %             \parhead{Combat Tactics} You are a master of all four elements, so you have an immense variety of options available to you - if you choose the right spells.
    %             You have a very high accuracy thanks to your Perception and a reasonably high power, so your primary role in combat will usually be to deploy the perfect damaging spell or debuff for the situation.
    %             Your skills and Elementalist abilities give you a lot of narrative power, so stay alert for opportunities to overcome challenges without needing to fight at all.

    %         \subsubsection{Nature Magic}
    %             \parhead{Species} Elf.
    %             % 0 3 0 0 3 2 base
    %             \parhead{Attributes} 0 Str, 3 Dex, -1 Con, 0 Int, 4 Per, 2 Wil (after species modifiers).
    %             \parhead{Class} Druid.
    %             \parhead{Archetypes} Nature Magic first, Nature Spell Mastery second, Elementalist third.
    %             \parhead{Insight Points} 1 point (unspent).
    %             \parhead{Skills} Awareness, Creature Handling, Knowledge (nature), Survival
    %             \parhead{Languages} Common, Elven, Halfling.
    %             \parhead{Equipment} Sickle, buckler, leather armor. As you gain levels, use the best light armor you can afford.
    %             \parhead{Mystic Spheres} Toxicology.
    %             \parhead{Suggested Feats} Sphere Focus: Verdamancy, Sphere Focus: Aquamancy, Herbalist
    %             \parhead{Combat Tactics} You are a master of plants and natural magic.
    %             Your spells excel at constraining and debilitating your foes, especially with poisons.
    %             Larger foes tend to be resistant to both poisons and movement impediments, so it's a good idea to have Reflex attacks like \spell{ensnaring grasp} and \spell{fire seeds} to deal with them.

    %         \subsubsection{Nature Spell Mastery}
    %             Use the typical character for the Nature Magic druid archetype.
    %             Even if you focus on spells through this archetype, you should generally still rank up your spells before improving your rank in this archetype.

    %         \subsubsection{Shifter}
    %             \parhead{Species} Human.
    %             % 3 2 2 1 0 0 base
    %             \parhead{Attributes} 3 Str, 3 Dex, 3 Con, 1 Int, 0 Per, 0 Wil.
    %             \parhead{Class} Druid.
    %             \parhead{Archetypes} Shifter first, Nature Magic second, Wildspeaker third.
    %             \parhead{Insight Points} 2 points.
    %             \parhead{Skills} Awareness, Poise, Athletics, Stealth, Survival
    %             \parhead{Languages} Common, Sylvan
    %             \parhead{Equipment} Natural weapon, buckler, chain shirt. As you gain levels, use the best light armor you can afford.
    %             \parhead{Mystic Sphere} Polymorph
    %             \parhead{Suggested Wild Aspects} Your choice of wild aspects has a significant effect on your capabilities, so choose wild aspects that match your goals.
    %             The Bear, Viper, and Wolf forms excel at dealing damage in combat.
    %             The Bull and Constrictor forms improve your ability to take unusual combat actions.
    %             Other forms can be useful in specific circumstances and out of combat.
    %             \parhead{Suggested Feats} Sphere Focus: Polymorph, Regenerator, Brawler, Savage
    %             \parhead{Combat Tactics} You are a lethal blend of claws and teeth.
    %             You can shift your form to gain the perfect abilities for your current circumstances, and your high physical attributes make you hard to kill and hard to ignore.
    %             Your flexibility between natural weapons, spells, and high physical skills give you a lot of options in and out of combat.
    %             In general, you do the most damage in close quarters where you can attack with your natural weapons, but you can use your spells to soften up strong enemies and finish off weakened enemies.

    %         \subsubsection{Wildspeaker}
    %             \parhead{Species} Gnome.
    %             \parhead{Attributes} -1 Str, 0 Dex, 4 Con, 0 Int, 3 Per, 3 Wil.
    %             \parhead{Class} Druid.
    %             \parhead{Archetypes} Wildspeaker first, Nature Magic second, Nature Spell Mastery third.
    %             \parhead{Insight Points} 1 point.
    %             \parhead{Skills} Awareness, Creature Handling, Knowledge (nature), Survival
    %             \parhead{Languages} Common, Gnome, Sylvan
    %             \parhead{Equipment} Sickle, buckler, scale mail. As you gain levels, use the best medium armor you can afford.
    %             \parhead{Mystic Sphere} Electromancy
    %             \parhead{Suggested Feats} Sphere Focus: Electromancy, Poise Specialization, Creature Handling Specialization, Toughness
    %             \parhead{Combat Tactics} You lead your faithful natural servant in battle.
    %             It distracts your enemies while you blast them with lightning from afar.
    %             Once you get a \textit{shrinking belt} or some other way to shrink yourself, you can ride your \textit{natural servant} into battle, which compensates for your short gnomish legs.
    %             If you are both lucky and persuasive, you be able to use your \textit{speak with animals} ability to convince an animal to aid you on your journey, at least for a short time, in addition to your \textit{natural servant}.

    %     \subsection{Fighter}

    %         \subsubsection{Combat Discipline}
    %             \parhead{Species} Dwarf.
    %             % 3 0 3 1 0 1 base
    %             \parhead{Attributes} 3 Str, -1 Dex, 4 Con, 1 Int, 0 Per, 2 Wil (after species modifiers).
    %             \parhead{Class} Fighter.
    %             \parhead{Archetypes} Combat Discipline first, Martial Mastery second, Sentinel third.
    %             \parhead{Insight Points} 2 points.
    %             \parhead{Skills} Awareness, Athletics, Endurance
    %             \parhead{Languages} Common, Dwarven, Orcish.
    %             \parhead{Equipment} Battleaxe, standard shield, scale mail. As you gain levels, buy a heavy crossbow and use the best heavy armor you can afford.
    %             You can switch between a shepherd's axe for hard to hit enemies, a battleaxe for multi-enemy fights or fights where you need the extra damage from holding it in two hands, and throwing axes when you need a ranged weapon.
    %             \parhead{Combat Styles} Rip and Tear, Unbreakable Defense
    %             \parhead{Suggested Feats} Shieldbearer, Toughness, Regenerator
    %             \parhead{Combat Tactics} You are extremely difficult to kill, and your ability to ignore and remove conditions makes it hard for your foes to whittle you down over time.
    %             You can charge confidently into the middle of battle, cutting down enemy ranged attackers regardless of their surrounding allies.
    %             Alternately, you can hold the line to protect your own allies.

    %         \subsubsection{Equipment Training}
    %             \parhead{Species} Halfling.
    %             % 0 2 2 1 2 1 base
    %             \parhead{Attributes} -1 Str, 3 Dex, 2 Con, 1 Int, 2 Per, 1 Wil.
    %             \parhead{Class} Fighter.
    %             \parhead{Archetypes} Equipment Training first, Martial Mastery second, Combat Discipline third.
    %             \parhead{Insight Points} 2 points.
    %             \parhead{Skills} Awareness, Poise, Flexibility, Stealth
    %             \parhead{Languages} Common, Gnomish, Halfling
    %             \parhead{Equipment} Kukri, standard shield, scale mail. As you gain levels, buy a longbow and use the best armor you can afford that allows you to apply your full Dexterity bonus to your Armor defense.
    %                 Keep an extra kukri with you so you can dual wield in fights where you don't need to use a shield.
    %             \parhead{Combat Styles} Flurry of Blows, Rip and Tear
    %             \parhead{Suggested Feats} Swiftrunner, Rapid Reaction, Executioner, Two-Weapon Fighting
    %             \parhead{Combat Tactics} You have an exceptionally high Armor defense, and your strikes are very accurate.
    %             However, your low Strength and small weapons mean that you don't deal a lot of damage.
    %             Focus on debilitating maneuvers like \maneuver{brow gash} or maneuvers that increase your damage like \maneuver{tear exposed flesh} to stay relevant in combat.
    %             When your Armor defense isn't as important, such as when fighting spellcasters, you can dual wield to increase your damage.
    %             Unlike most fighters, you are very agile and stealthy, so you can accompany rogues on scouting missions.

    %         \subsubsection{Martial Mastery}
    %             \parhead{Species} Human.
    %             % 3 0 2 1 2 0 base
    %             \parhead{Attributes} 3 Str, 0 Dex, 3 Con, 1 Int, 3 Per, 0 Wil.
    %             \parhead{Class} Fighter.
    %             \parhead{Archetypes} Martial Mastery first, Combat Discipline second, Tactician third.
    %             \parhead{Insight Points} 2 points.
    %             \parhead{Skills} Awareness, Athletics, Endurance
    %             \parhead{Languages} Common, Giantish, Orcish.
    %             \parhead{Equipment} Broadsword, standard shield, scale mail. As you gain levels, buy throwing axes and use the best heavy armor you can afford.
    %             \parhead{Combat Styles} Ebb and Flow, Herald of War, Rip and Tear
    %             \parhead{Suggested Feats} Executioner, Blindfighter, Toughness
    %             \parhead{Combat Tactics} You have great versatility in combat.
    %             You have a large number of manuvers, and by mixing in thrown weapons and shouts, you can attack at multiple ranges and hit multiple defenses.
    %             When your shield is unnecessary, you can hold your broadsword in two hands to improve your already respectable damage.
    %             Many of your maneuvers work at any range, so you aren't forced to fight in melee against highly mobile or excessively lethal foes.
    %             In addition to using maneuvers, you can coordinate your allies with battle tactics.

    %         \subsubsection{Sentinel}
    %             \parhead{Species} Dwarf.
    %             % 3 0 2 0 3 0 base
    %             \parhead{Attributes} 3 Str, -1 Dex, 3 Con, 0 Int, 3 Per, 1 Wil.
    %             \parhead{Class} Fighter.
    %             \parhead{Archetypes} Sentinel first, Martial Mastery second, Equipment Training third.
    %             \parhead{Insight Points} 1 point.
    %             \parhead{Skills} Awareness, Endurance, Intimidate
    %             \parhead{Languages} Common, Dwarven, Orcish
    %             \parhead{Equipment} Shepherd's axe, standard shield, scale mail. As you gain levels, buy a throwing axes, a heavy crossbow, and a greataxe and use the best heavy armor you can afford.
    %             \parhead{Combat Styles} Unbreakable Defense
    %             \parhead{Suggested Feats} Shieldbearer, Toughness, Regenerator
    %             \parhead{Combat Tactics} You hold the line in the middle of the fray, protecting your allies all over the battlefield.
    %             You have an unmatched ability to constrain your foes' movement and force them to pay attention to you, limiting their ability to harm your allies.
    %             Your damage is reasonable, but your main focus should be on defense so you can protect your more vulnerable and reckless allies.
    %             A shepherd's axe is convenient because it allows you to attack foes who try to keep their distance from you.
    %             However, when you need to deal damage, consider switching to a greataxe or dual-wielding a second shepherd's axe instead of holding a shield.

    %         \subsubsection{Tactician}
    %             \parhead{Species} Human.
    %             % 2 2 1 1 2 0 base
    %             \parhead{Attributes} 2 Str, 3 Dex, 1 Con, 2 Int, 2 Per, 0 Wil.
    %             \parhead{Class} Fighter.
    %             \parhead{Archetypes} Tactician first, Martial Mastery second, Equipment Training third.
    %             \parhead{Insight Points} 2 points (1 for Mobile Hunter combat style, 1 unspent).
    %             % 3 fighter + 2 int + 1 human = 6 trained
    %             \parhead{Skills} Awareness, Deduction, Endurance, Knowledge (dungeoneering, local), Medicine
    %             \parhead{Languages} Common, Draconic, Elven
    %             \parhead{Equipment} Longhammer, standard shield, chain shirt. As you gain levels, buy a longbow and use the best armor you can afford that allows you to apply your full Dexterity bonus to your Armor defense.
    %             \parhead{Combat Styles} Brute Force, Mobile Hunter
    %             \parhead{Suggested Feats} Precognition, Leadership, Medicine Specialization
    %             \parhead{Combat Tactics} You have a wealth of options in combat.
    %             You can buff your allies' attacks, defend your allies, deal damage, or debuff your foes.
    %             Use whichever battle tactics are most relevant to the current situation.
    %             Your longhammer and high mobility allows you to keep your distance in combat while knocking your foes into more tactically advantageous positions.
    %             Since you don't use a shield, you don't have the defensive power of most fighters, so you can't just charge heedlessly into the fray.

    %     \subsection{Monk}

    %         \subsubsection{Airdancer}
    %             \parhead{Species} Elf.
    %             % 2 2 0 0 2 2 base
    %             \parhead{Attributes} 2 Str, 3 Dex, -1 Con, 0 Int, 2 Per, 2 Wil.
    %             \parhead{Class} Monk.
    %             \parhead{Archetypes} Airdancer first, Esoteric Warrior second, Ki third.
    %             \parhead{Insight Points} 1 point.
    %             \parhead{Skills} Poise, Athletics, Flexibility, Stealth
    %             \parhead{Languages} Common, Elven, Halfling
    %             \parhead{Equipment} Two jitte. Use your \textit{ki barrier} for your body armor.
    %             \parhead{Combat Styles} Perfect Precision
    %             \parhead{Suggested Feats} Poise Specialization, Swiftrunner, Two-Weapon Fighting
    %             \parhead{Combat Tactics} You are highly acrobatic in combat, leaping around your opponents with ease.
    %             Once you can jump over enemies, you can start ignoring attempts to block your movement.
    %             The Leap of the Heavens \textit{ki manifestation} can help you reach that point quickly.
    %             You are highly accurate, and your high Dexterity helps both your defenses and your skills.
    %             If you are in physical danger, you can sheathe one of your kamas and use your \textit{ki barrier} as a shield to further increase your Armor defense.
    %             However, your middling Constitution means you should pick your fights carefully.
    %             Use your mobility to pick your battles and avoid being surrounded unnecessarily.

    %         \subsubsection{Esoteric Warrior}
    %             \parhead{Species} Human
    %             % 2 2 2 1 1 0 base
    %             \parhead{Attributes} 2 Str, 3 Dex, 2 Con, 1 Int, 2 Per, 0 Wil.
    %             \parhead{Class} Monk.
    %             \parhead{Archetypes} Esoteric Warrior first, Perfected Form second, Transcendent Sage third.
    %             \parhead{Insight Points} 2 points (1 for Flurry of Blows combat style, 1 point unspent).
    %             \parhead{Skills} Awareness, Poise, Athletics, Flexibility, Endurance, Stealth
    %             \parhead{Languages} Common, Draconic, Elven
    %             \parhead{Equipment} Two kunai, chain shirt. As you gain levels, buy spare kunai, and use the best light armor you can afford.
    %             \parhead{Combat Styles} Dirty Fighting, Flurry of Blows
    %             \parhead{Suggested Feats} Brawler, Juggernaut, Swiftrunner, Two-Weapon Fighting
    %             \parhead{Combat Tactics} You can beat your opponents to death with nothing more than your bare hands.
    %             Your primary combat strategy is generally to grapple, trip, or otherwise debuff your opponents with your free hands before you pummel them into submission.
    %             You have a high movement speed, and you can take advantage of that by rushing down enemies who would prefer to keep their distance.
    %             When that is combined with your immunity to many common debuffs, you are exceptionally effective against enemy spellcasters.
    %             If you find yourself fighting more martially skilled foes, you may need to keep your distance with kunai or a bow, at least until they are weakened.

    %         \subsubsection{Ki}
    %             \parhead{Species} Halfling
    %             % 0 2 3 1 0 2 base
    %             \parhead{Attributes} -1 Str, 2 Dex, 3 Con, 1 Int, 0 Per, 3 Wil.
    %             \parhead{Class} Monk.
    %             \parhead{Archetypes} Ki first, Esoteric Warrior second, Transcendent Sage third.
    %             \parhead{Insight Points} 2 points.
    %             \parhead{Skills} Poise, Flexibility, Endurance, Knowledge (arcana), Survival, Stealth
    %             \parhead{Languages} Common, Draconic, Elven
    %             \parhead{Equipment} Two kama. Use your \textit{ki barrier} for your body armor. As you gain levels, buy spare kunai.
    %             \parhead{Combat Styles} Flurry of Blows, Mobile Hunter
    %             \parhead{Suggested Feats} Two-Weapon Fighting, Ghostblade, Iron Will, Spellwarped
    %             \parhead{Combat Tactics} Although you appear small and physically weak, your attacks hit hard thanks to your \textit{ki energy} ability.
    %             You can use a variety of ki manifestations to have surprising effects in combat.
    %             Look for tricky combinations, like tripping your foes at a distance with \ability{extend the flow of ki} or using \ability{burst of blinding speed} to increase the power of movement-based maneuvers.
    %             You can also use your ki manifestations to augment your skills in non-combat situations.
    %             Your defenses are high and well-rounded, and you have immunities to a variety of common debuffs, so you can fight aggressively in combat.
    %             Generally, you should dual-wield kama, but you can drop to a single kama if you need more Armor defense, or you can switch to throwing kunai to hit distant foes.

    %         \subsubsection{Perfected Form}
    %             \parhead{Species} Human
    %             % 3 2 2 1 0 0 base
    %             \parhead{Attributes} 3 Str, 3 Dex, 3 Con, 1 Int, 0 Per, 0 Wil.
    %             \parhead{Class} Monk.
    %             \parhead{Archetypes} Perfected Form first, Esoteric Warrior second, Airdancer third.
    %             \parhead{Insight Points} 2 points (1 for Mobile Hunter combat style, 1 point unspent).
    %             \parhead{Skills} Poise, Athletics, Flexibility, Endurance, Stealth
    %             \parhead{Languages} Common, Giantish, Orcish
    %             \parhead{Equipment} Two kunai, chain shirt. As you gain levels, buy spare kunai, and use the best light armor you can afford.
    %             \parhead{Combat Styles} Dirty Fighting, Mobile Hunter
    %             \parhead{Suggested Feats} Brawler, Juggernaut, Swiftrunner, Two-Weapon Fighting
    %             \parhead{Combat Tactics} Your general fighting style is the same as the Esoteric Warrior sample character.
    %             Your main differentiating factor is that you have even greater martial aptitude and durability.
    %             However, you are less able to avoid to debilitating conditions, so be careful when fighting magical foes.

    %         \subsubsection{Transcendent Sage}
    %             The transcendent sage archetype by itself does not strongly influence your character's fighting style or abilities.
    %             It provides a variety of passive abilities and immunities that require other archetypes to make a compelling character concept.
    %             A typical character focusing on this archetype would be similar to the Ki character.
    %             If you want to be more martially inclined, follow the Esoteric Warrior character.

    %     \subsection{Paladin}

    %         \subsubsection{Devoted Paragon}
    %             The devoted paragon archetype can have different play styles based on your devoted alignment.
    %             The character below is reasonable for any alignment other than evil.
    %             An evil devoted paragon might focus more on inflicting debilitating conditions with spells from the \sphere{enchantment} or \sphere{vivimancy} \glossterm{mystic spheres}.

    %             \parhead{Species} Dwarf
    %             % 2 0 3 0 2 1 base
    %             \parhead{Attributes} 2 Str, -1 Dex, 4 Con, 0 Int, 2 Per, 2 Wil.
    %             \parhead{Class} Paladin.
    %             \parhead{Archetypes} Devoted Paragon first, Divine Magic second, Stalwart Guardian third.
    %             \parhead{Insight Points} 1 point.
    %             \parhead{Skills} Endurance, Persuasion, Social Insight
    %             \parhead{Languages} Common, Dwarven, Giantish
    %             \parhead{Equipment} Morning star, standard shield, scale mail. As you gain levels, buy a heavy crossbow and use the best heavy armor you can afford.
    %             \parhead{Mystic Sphere} Prayer.
    %             \parhead{Suggested Feats} Leadership, Weapon Focus, Shieldbearer, Sphere Focus: Prayer
    %             \parhead{Combat Tactics} You are a beacon to guide your allies in combat.
    %             You should stay close to the front lines to ensure that your allies gain the benefit of your \textit{aligned aura} ability and take advantage of your reasonable melee damage.
    %             You are very durable, and you can blanket your allies with powerful buffs from the \sphere{prayer} mystic sphere.
    %             At high levels, make sure you have a special strike ability like \spell{exalted strike} or a \mitem{powerstrike} weapon.

    %         \subsubsection{Divine Magic}

    %             \parhead{Species} Human
    %             % 1 0 2 2 1 2 base
    %             \parhead{Attributes} 1 Str, 0 Dex, 2 Con, 3 Int, 1 Per, 3 Wil.
    %             \parhead{Class} Paladin.
    %             \parhead{Archetypes} Divine Magic first, Divine Spell Expertise second, Devoted Paragon third.
    %             \parhead{Insight Points} 3 points (2 points for an additional mystic sphere (Vivimancy), 1 point for an additional rank 1 spell).
    %             \parhead{Skills} Awareness, Deduction, Endurance, Knowledge (local, religion), Medicine, Persuasion
    %             \parhead{Languages} Common, Dwarven, Giantish
    %             \parhead{Equipment} Battleaxe, standard shield, scale mail. As you gain levels, use the best heavy armor you can afford and switch to using an implement instead of a weapon.
    %             \parhead{Mystic Sphere} Channel Divinity, Vivimancy.
    %             \parhead{Suggested Feats} Sphere Focus: Channel Divinity, Sphere Focus: Vivimancy, Boongiver
    %             \parhead{Combat Tactics} Paladins make unusual spellcasters.
    %             They lack the flexibility of more dedicated spellcasting classes because they lack the \ability{metamagic} ability and access to a variety of mystic spheres.
    %             However, your \textit{divine spell versatility} and \textit{divine conduit} give you a unique combat style that rewards staying directly on the front lines.
    %             In addition, you have an unusually high power thanks to \textit{wellspring of power} and \textit{paragon power}, so your damage and healing spells can be quite potent.
    %             Make sure you have single-target damage spells like \spell{mystic bolt} and \spell{lifesteal} to make the most of your abilities.

    %         \subsubsection{Divine Spell Expertise}
    %             Use the typical character for the Divine Magic paladin archetype.
    %             Even if you focus on spells through this archetype, you should generally still rank up your spells before improving your rank in this archetype.

    %         \subsubsection{Stalwart Guardian}

    %             \parhead{Species} Halfling
    %             % 0 3 3 0 0 2 base
    %             \parhead{Attributes} -1 Str, 3 Dex, 3 Con, 0 Int, 0 Per, 3 Wil (after species modifiers).
    %             \parhead{Class} Paladin.
    %             \parhead{Archetypes} Stalwart Guardian first, Divine Magic second, Zealous Warrior third.
    %             \parhead{Insight Points} 1 point.
    %             \parhead{Skills} Endurance, Medicine, Poise
    %             \parhead{Languages} Common, Dwarven, Giantish
    %             \parhead{Equipment} Smallsword, standard shield, scale mail. As you gain levels, buy spare throwing axes and use the best light armor you can afford.
    %             \parhead{Mystic Sphere} Fabrication.
    %             \parhead{Suggested Feats} Shieldbearer, Sphere Focus: Prayer
    %             \parhead{Combat Tactics} You have exceptionally high defenses, and you can heal your less well protected allies.
    %             Your damage is reasonable thanks to your \ability{smite} ability, which allows you to use your high Willpower in place of your poor Strength.
    %             You can use spells like \spell{mystic barrier} and \spell{blade barrier} to control the battlefield and funnel your foes into positions where you can make the most of your defensive abilities.

    %         \subsubsection{Zealous Warrior}

    %             \parhead{Species} Elf
    %             % 3 0 0 0 2 3 base
    %             \parhead{Attributes} 3 Str, 0 Dex, -1 Con, 0 Int, 3 Per, 3 Wil (after species modifiers).
    %             \parhead{Class} Paladin.
    %             \parhead{Archetypes} Zealous Warrior first, Divine Magic second, Stalwart Guardian third.
    %             \parhead{Insight Points} 1 point.
    %             \parhead{Skills} Awareness, Persuasion, Social Insight
    %             \parhead{Languages} Common, Dwarven, Giantish
    %             \parhead{Equipment} Greatmace, scale mail. As you gain levels, buy a longbow and use the best heavy armor you can afford.
    %             \parhead{Mystic Sphere} Channel Divinity.
    %             \parhead{Suggested Feats} Greatweapon Warrior, Spellsword, Celestial Ancestry
    %             \parhead{Combat Tactics} You are an extremely dangerous martial combatant.
    %             Thanks to your \ability{smite} ability, you have high damage and devastatingly powerful critical hits.
    %             Your defenses are unimpressive, but you have enough durability from heavy armor to avoid being too squishy.
    %             The biggest weakness you have is your slow speed and vulnerability to debilitating conditions.
    %             Make sure you take some ranged spells like \spell{retributive judgment} to deal with foes who try to avoid your greatmace.

    %             Alternately, you could focus more heavily on ranged combat, since \ability{smite} can be used with \weapontag{Projectile} weapons.
    %             If that is your goal, consider the Mystic Archer and Sniper feats instead of Greatweapon Warrior and Spellsword.
